% Fragmento científico exclusivo de 01_electron.tex, líneas 120--197.
% El prefijo 1--118 coincide byte a byte con la ontología principal ya
% publicada; este fragmento conserva únicamente la formulación cuántica
% adicional sin repetir definiciones idénticas.
\section{Ontologie quantique de la particule}
\Needspace{7\baselineskip}
\subsection{État, histoire et réalisation}
Un état généalogique est une section compatible du système d'extensions du TPK. Une histoire observable est la projection chronologique de cette section sur un lecteur. Une réalisation physique est la représentation de l'histoire dans une droite dimensionnelle, un espace de Hilbert, une algèbre de champs ou un codomaine topologique. Ces trois notions ne sont pas interchangeables. L'état porte l'identité ; l'histoire porte la preuve ; la réalisation porte l'observable.

Soit \(\mathcal X_{\rm TPK}^{\rm enr}\) l'espace des états enrichis et \(p\) le lecteur observable. Une particule est une classe \([X]\) de sections persistantes telle que

\[
p(X_{K+1})|_K=p(X_K),
\]
la loi de transition conserve les relations de l'enregistrement généalogique et que le prolongement ne se termine pas sur un horizon d'obstruction. Le parcours

\[
\gamma_X=(p(X_0),p(X_1),p(X_2),\ldots)
\]
est l'histoire visible de la particule. Il n'est pas, par définition, une trajectoire spatiale. Une réalisation métrique ultérieure peut le convertir en courbe, orbite ou propagateur lorsqu'elle construit l'entrelaceur correspondant.

L'identité physique se conserve sous raffinement lorsque les carrés

\[
\begin{CD}
X_{K+1} @>{F_{K+1}}>> X'_{K+1}\\
@V{r_{K+1,K}}VV @VV{r'_{K+1,K}}V\\
X_K @>>{F_K}> X'_K
\end{CD}
\]
commutent. La masse, la charge et le spin sont des lecteurs distincts de ce même antécédent. L'égalité de l'un n'impose pas l'égalité des autres.

\Needspace{7\baselineskip}
\subsection{Représentation quantique de la particule comme section locale persistante de l'état généalogique}
Une particule HMT--MD est une section persistante d'une extension enrichie qui satisfait quatre conditions :

\begin{enumerate}
\item elle possède un parcours observable compatible avec la monodromie nonadique ;
\item elle conserve un enregistrement généalogique sous raffinement ;
\item elle admet une représentation d'orientation, de charge, de spin et de statistique ;
\item elle présente un résultat spectral stable dans le codomaine physique déclaré.
\end{enumerate}
La stabilité ne signifie pas l'immobilité. Elle signifie que les changements de phase, de feuillet et de mémoire appartiennent à une même classe généalogique. La particule peut clore sa phase sans clore son état, changer de feuillet tout en conservant sa masse, ou traverser plusieurs cellules en maintenant la même identité d'extension.

Une espèce est une famille de particules partageant le classificateur interne

\[
\mathfrak s=
(\mathfrak q,\mathfrak j,\mathfrak c,\mathfrak f,\mathfrak g,
\mathfrak r,\mathfrak o),
\]
où \(\mathfrak q\) est la représentation de charge, \(\mathfrak j\) l'holonomie de spin, \(\mathfrak c\) l'orbite de couleur, \(\mathfrak f\) la saveur, \(\mathfrak g\) la génération, \(\mathfrak r\) la famille de parcours et \(\mathfrak o\) l'opérateur de masse. Une espèce non centrale peut correspondre à une multisection complète ; exiger d'elle un point APP unique avant de déclarer le couplage physique constitue une erreur de type.

\Needspace{7\baselineskip}
\subsection{Conjugaison quantique de l'antiparticule par inversion de feuillet d'une section persistante}
L'involution dodécaphasique est

\[
C_M(\tau)=-\operatorname{rev}(\tau).
\]
Elle agit sur l'histoire orientée et se prolonge à l'enregistrement généalogique. La charge et l'orientation linéaire sont impaires ; la mémoire quadratique est paire. Si \(q\) est le lecteur de charge,

\[
q(C_MX)=-q(X).
\]
L'égalité des masses exige en outre la covariance opératorielle. Si \(K\) réalise la conjugaison dans l'espace physique et

\[
K\widehat M_XK^{-1}=\widehat M_{\bar X},
\]
en transportant aussi les canaux et la prescription causale lorsque \(K\) est antiunitaire, alors \(m_{\bar X}=m_X\). La seule parité de l'enregistrement ne remplace pas cette prémisse.

Particule et antiparticule sont deux sections orientées d'une même bande généalogique. Dans la réalisation de Möbius, un tour échange le feuillet et deux tours rétablissent l'orientation. La conclusion ne procède pas du dessin d'un ruban tordu : elle exige l'involution, le relèvement au revêtement double, la boucle de retour et l'holonomie de signe.

Pour l'enregistrement électronique \(\tau_e=+++---+++---\),

\[
C_M\tau_e=\tau_e.
\]
Le mot central est fixe, tandis que la représentation de charge distingue électron et positron. On obtient ainsi une même masse sans doubler le parcours temporel. La clôture à 216 pas ne signifie pas 108 pas de particule et 108 d'antiparticule : elle signifie le rétablissement du relèvement cinématique orienté.

\Needspace{7\baselineskip}
